
La arquitectura distribuida de \textit{NodeAlert IoT} sienta las bases para diversas líneas de expansión tecnológica. Se plantean las siguientes propuestas para evolucionar el sistema hacia un entorno de monitoreo inteligente y de escala industrial:

\subsection{Inteligencia Artificial y Análisis Predictivo}
Una evolución natural del proyecto consiste en la implementación de modelos de \textit{Machine Learning} para la predicción de incendios. En lugar de depender de umbrales fijos, el sistema podría analizar patrones históricos de temperatura y humedad para anticipar condiciones de riesgo antes de que ocurran. Según \citet{deitel2004}, la integración de algoritmos de análisis de datos permite transformar la información cruda en conocimiento accionable, mejorando la proactividad del sistema.

\subsection{Optimización de Protocolos y Seguridad}
Se proyecta la implementación de estándares de seguridad más robustos, como TLS/SSL para las comunicaciones MQTT, asegurando la privacidad de los datos en entornos corporativos. Asimismo, se explorará la optimización de los registros de comunicación a nivel de bit, basándose en las técnicas de programación de periféricos de \citet{mazidi2014}, para reducir aún más el consumo energético en nodos alimentados por energía solar.